\section{Further research questions and proposed work programme}
\label{sec:research}

The next phase should be driven by precise missing lemmas and adversarial
families. A successful experiment is useful when it identifies a parameter or
falsifies an attractive universal claim. The figure-eight improvement and the
Fibonacci-word refusal together give a concrete starting point: signed
substitution matters, but singleton access is not a complete basis-change
mechanism. The following questions are ordered roughly from immediate native
work to the global geometric bottlenecks.

\subsection{Immediate terminal and implementation questions}

\begin{question}[Recover the two admitted disk-base stalls]
Can explicit source-aware Nielsen transformations expose defining generators
in the two remaining disk-base word stalls, while keeping a polynomial total
word-size budget? First retain their full signed presentations and compare
different generator bases. Require exact automorphism certificates and replay
of the original words. Success means new whole-cut certificates or a small,
repeatable obstruction to the proposed basis policy; a change in abelianized
rank alone does not count.
\end{question}

\begin{question}[A stronger bounded algebraic policy]
What is the smallest additional move family that resolves automorphic
images such as $\langle a,b\mid\varphi^k(a),\varphi^k(b)\rangle$ efficiently
without making discovery exponential in the rank? Candidate moves include
certified Nielsen basis changes, relator multiplication and controlled
temporary length growth. Charge both number of moves and intermediate literal
or compressed bit size. A short certificate found by unbounded search is a
different achievement from a bounded discovery algorithm.
\end{question}

\begin{question}[Package the early homology obstruction]
Can the partial-trace condition $z>h_3$ be exported as a compact independent
whole-cut obstruction certificate? The theorem already proves its soundness.
The implementation task is to retain the verified prefix, authenticate the
source count and expose a clearly typed result. Such a result can prune a
terminal candidate while preserving the possibility that a different normal
surface gives a useful hierarchy.
\end{question}

\begin{question}[Reuse authenticated source work across the portfolio]
Can the geometric and group terminals share one immutable source/boundary
packet and one explicit work ledger? The present opt-in wrapper repeats source
preparation after a refusal. A useful refactor would reduce measured failed-path
latency and eliminate the two-stage orbit allowance, while preserving legacy
certificate output and callback semantics. Measure complete successful queries
and inconclusive queries separately, including packet serialization.
\end{question}

\begin{question}[Incremental relator substitution]
Can inverted occurrence lists, lazy substitutions or straight-line programs
reduce the current $O(N^4)$ discovery bound and its constants without changing
the selected minimum-length move? An occurrence index can avoid scanning words
that do not contain the pivot, but exact cyclic cancellation and candidate
comparison still need a proof. Record total changed-word work, transient word
length and certificate replay cost, rather than only counting chosen pivots.
\end{question}

\subsection{Larger source classes and geometric certificates}

\begin{question}[Extend the compiler beyond disk bases]
For a compact base surface with boundary, its group is free but its boundary
curves have specified words. Can a graph-of-spaces compiler retain these words,
cyclic edge groups for full annuli, and trivial edge groups for intervals with
an active size controlled by ports and encoded base complexity? A boundary loop
of a positive-genus one-boundary surface is not a free generator. Replacing it
by one would be unsound. Very large encoded genus may also force a long
commutator word unless a justified compressed representation is used.
\end{question}

\begin{question}[Terminal residual kernels in real hierarchies]
After geometric peeling and signed simplification, is the terminal residual
size logarithmic in the original diagram size on a geometrically defined
class? The exact semantic kernel in the preceding section makes this a
well-posed parameter question. The old linear-deferral ball family no longer
obstructs it, since signed elimination solves every block. The Fibonacci basis
family shows that raw singleton residual size is sensitive to coordinates.
A useful theorem must therefore specify allowed source-certified basis changes.
\end{question}

\begin{question}[Recover source-coordinate meridians]
When the Euler/free-group terminal certifies $H_g$, can a separate procedure
construct a complete meridian system with bounded compressed descriptions and
verified boundary carriers? The current proof uses the Loop Theorem and
Poincar\'e only for existence. Converting an abstract group basis into disks is
an additional geometric problem. A first milestone is the explicit
figure-eight genus-two component, with independently checked embedded disks
and transports to the original source charts.
\end{question}

\begin{question}[Essential boundary patterns at an algebraic terminal]
Can a bounded marked pattern be checked for the properties needed by a
hierarchy once the ambient piece is certified to be a handlebody? Merely
knowing the genus or the abstract group leaves the pattern's embedding unknown.
The desired packet should retain cyclic order, side information, slopes where
relevant, and witnesses that rule out violating disks. This is a natural point
to connect finite boundary graphs with source-coordinate meridian certificates.
\end{question}

\begin{question}[Turn apparent basis obstructions into geometric families]
Can the algebraically incomplete Fibonacci-word examples be realized and
exported by the native normal-cut constructor with controlled triangulation
size? Conversely, do the native disk-annulus sources satisfy an extra
restriction that rules out some arbitrary finite-presentation obstructions?
Producing actual face pairings and normal coordinates, rather than a plausible
attaching-word diagram, is the required experiment.
\end{question}

\subsection{The global quasi-polynomial bottlenecks}

\begin{question}[Find an amortized geometric branch potential]
Can a compressed hierarchy admit a nonnegative potential whose decrease pays
for $\log b$ at a node with $b$ genuine successor alternatives? Polynomially
many deterministic repairs are harmless if they do not replenish this
potential. Candidate ingredients include boundary pattern complexity,
nonproduct core size, a hierarchy surface's essential topology and the residual
group kernel. The local inequality in the branch-budget theorem is a concrete
target; mere lexicographic termination does not imply it.
\end{question}

\begin{question}[Construct a small marked interface quotient]
Is there a source-liftable signature on $O(\log n)$ marked interface records
that preserves every relevant continuation and terminal condition? The
signature must account for mapping-class information that group rank,
component incidence and Euler characteristic forget. For every proposed
signature, actively search for two equal-key inputs with different outcomes
under the same allowed continuation. A collision of this kind is more
informative than a high cache-hit rate on easy inputs.
\end{question}

\begin{question}[Bound candidate surfaces, not just one chosen surface]
Can the algorithm generate a complete family of geometrically useful normal
cuts in polynomial or quasi-polynomial total work at each essential branching
stage? Efficiently finding one nonseparating surface does not show it is enough
for every later repair. The statement should specify which alternatives cover
all successful hierarchies and charge the cost of rejected candidates,
coordinate generation and source validation. A short existential hierarchy
certificate is not a deterministic search bound.
\end{question}

\begin{question}[Control marking growth under repeated cuts]
Can every successor be represented in size polynomial in the original diagram
length, including all coordinates, interval endpoints, boundary patterns and
transport proofs? A core with $O(n)$ handles can still carry enormous boundary
word data. The bound must cover the complete state and certificate, not only
the geometric core or the new active group presentation.
\end{question}

\begin{question}[Connect terminal improvements to a diagram-level schedule]
Which expensive current recognizer branches can be avoided when a supplied cut
receives the new certificate? Implement a bounded experimental schedule that
records diagram-to-exterior provenance, candidate-surface cost, source cuts,
all retries and final independent verification. Report added successful knot
verdicts and total latency, preserving timeout and inconclusive semantics. The
current three added cut successes alone do not predict that outcome.
\end{question}

\begin{question}[Formalize the smallest trusted mathematical interface]
Can the disk-base graph-of-spaces compiler, signed Tietze trace and Euler count
ledger be formalized separately from a full normal-surface library? The finite
combinatorial layer is a promising first target. The topological theorem could
initially be stated with explicit classical assumptions, then connected to a
formal 3-manifold development. This would expose exactly which source geometry
invariants must be certified rather than inferred from implementation history.
\end{question}

\subsection{Suggested milestones}

The first milestone is a reviewed integration of the present opt-in terminal,
with the three added native cases and five Euler-obstruction packets kept as
permanent regressions. Next, extend the residual algebraic policy and measure
its effect on actual source presentations; do not optimize arbitrary group
examples without confirming that they model the geometric bottleneck. In
parallel, construct source-coordinate marked meridian data for the figure-eight
target and evaluate the nondisk graph-of-spaces compiler.

The decisive theoretical milestone is then either a logarithmic amortized
branch budget or a complete small marked-state quotient, with polynomial state
size and proof lifting. The new terminal makes one local operation finite and
cheap in the relevant encoding. It also supplies exact refusal information.
Those are useful components of a complexity proof, while the global geometric
claims above remain the substantive research problem.
